\section{Settings}
\label{settings}

\subsection{\hr{Segelflugzeug Polare}}

\tr{Um die korrekten Sollfahrt Informationen zu erhalten, müssen Sie die richtigen Polarenwerte für
Ihren Segelflugzeugtyp einstellen. Das LARUS Vario Display verfügt werksseitig über mehr als 200
Polaren von verschiedenen Segelflugzeugen. Wählen Sie die richtige oder nächstliegende Polare unter
aus. Nach der Auswahl des Segelflugzeugtyps sollten Sie die Leergewichtsdaten Ihres Segelflugzeugs
aktualisieren.

Sollten Sie Ihren Segelflugzeugtyp nicht in der Liste finden, können Sie eine beliebige
Segelflugzeug-Polare auswählen und die einzelnen Einstellungen auf die Werte Ihrer
Segelflugzeug-Polare ändern. Der Name des Segelflugzeugs kann nicht geändert werden.}


\picnote
    {pictograph-yellow-warning.pdf} {\tr{Bei jedem Aufrufen der Segelflugzeugliste in den
    Polareinstellungen werden alle zuvor geänderten Werte durch die fest programmierten Werte des
    ausgewählten Segelflugzeugs ersetzt.}}
\picnote
    {pictograph-yellow-warning.pdf} {\tr{Achten Sie darauf, dass der maximale Wasserballast mit den
    Angaben von XCSoar übereinstimmt, da ansonsten der Abgleich des Wasserballastes mit XSCoar nicht
    korrekt funktoniert.}}


\subsection{\hr{Kalibrierung der Larus Senor Box}}

\tr{Bevor Sie Ihren ersten Flug beginnen, müssen die Lagesensoren des LARUS-Sensoreinheit präzise
eingestellt werden. Die Kalibrierungsschritte durch einen einfachen Ablauf durchgeführt, der
nachfolgend beschrieben ist und über Funktionen im LARUS Vario Display angestoßen werden.
Die Kalibrierung erfolgt in zwei Abschnitten.}

\subsubsection{\hr{Erste Kalibrierung am Boden:}}

\tr{Bauen Sie Ihren Segelflugzeug auf und stellen Sie ihn auf eine ebene Fläche. Nachdem Sie die
einzelnen Positionen eingenommen haben, warten Sie, bis keine Vibrationen mehr im Flugzeug zu spüren
sind, bevor Sie mit der Kalibrierung fortfahren. Verwenden Sie keinen Heckwagen, um die vertikale
Achse Ihres Segelflugzeugs während der folgenden Vorgänge zu fixieren.}

\begin{itemize}
    \item   \tr{Legen Sie den linken Flügel ab, warten Sie kurz und rufen Sie die Funktion
            Einstellungen / Sensoreinheit / Linker Flügel nach unten auf.}
    \item   \tr{Legen Sie den rechten Flügel ab, warten Sie kurz und rufen Sie die Funktion
            Einstellungen / Sensoreinheit / Rechter Flügel nach unten auf.}
    \item   \tr{Halten Sie den Flügel horizontal, warten Sie, rufen Sie die Funktion „Einstellungen
            / Sensoreinheit / Flügel gerade“ auf.}
    \item   \tr{Für diesen Schritt ist es wichtig, alle drei zuvor genannten Schritte durchzuführen.
            Die Reihenfolge der Schritte spielt keine Rolle, jedoch müssen sie vollständig
            abgeschlossen sein. Rufen Sie die Funktion Einstellungen / Sensoreinheit / Ausrichtung
            berechnen auf.}
\end{itemize}


\subsubsection{\hr{Feinjustage in der Luft:}}

\tr{Die Neigungswinkelkalibrierung wird während des Fluges durchgeführt. Es wird empfohlen, diesen
 Schritt in einem Flug durchzuführen, der nicht durch thermische Böen gestört wird. Richten Sie
 Ihren Segelflugzeug bei der Geschwindigkeit mit der besten Gleitzahl aus (wenn Sie Wölbklappen
 haben, stellen Sie diese auf diese Geschwindigkeit ein). Rufen Sie „Einstellungen / Sensoreinheit /
 Geradeausflug“ auf. Damit sind Sie fertig. Sie können die Kalibrierung überprüfen, indem Sie zur
 AHRS-Anzeige wechseln.}
